\chapter{The Response Model and Anti-Circularity}\label{ch06:response}

The center of CAS-Ohm 1.0 is one small equation and one large refusal. The
equation is the candidate prediction
\[
\hat y \;=\; \frac{a}{r_{\mathrm{eff}}},
\]
a proposed local linear relation between an intervention dose and an
observed response under a fixed calibration profile. The refusal is the
anti-circularity rule: the candidate rejects, categorically and at the type
level, the procedure that computes $r = a/y$ on a sample, computes
$\hat y = a/r$ on the same sample, observes $\hat y = y$, and announces a
validated law. This chapter develops both---the equation with its
admissibility conditions, and the refusal with its statistical
justification---then works through hypothetical numeric illustrations,
distinguishes verification from validation in Sargent's sense
\cite{sargent}, and closes by delimiting exactly what a validated response
model would and would not license.

The claim posture is stated once and holds throughout: everything here is
\textsc{ring 4 / not licensed / held}. The response model is a registered
hypothesis with an implemented guard apparatus, not an established law; no
live calibration exists; every number in the worked examples is fabricated
for exposition and is labeled as such where it appears.

\section{The Candidate Prediction}\label{ch06:prediction}

\begin{definition}[Candidate response prediction]\label{ch06:def:prediction}
Fix a calibration profile $P$ (Chapter~\ref{ch05:profile}) and a registered
test point $t$ of $P$. Let $a$ be an alpha dose (quantity kind
\texttt{DRIVE\_INTERVENTION}, unit \texttt{ALPHA}) within $P$'s declared
alpha range, and let $r_{\mathrm{eff}}$ be an effective resistance at $t$
(quantity kind \texttt{EFFECTIVE\_RESISTANCE}, unit
\texttt{ALPHA\_PER\_ACTIVATION}, strictly positive). The candidate
prediction for the directional activation response is
$\hat y = a / r_{\mathrm{eff}}$, a quantity in \texttt{ACTIVATION} units,
asserted---as hypothesis, not fact---to approximate the observed response
$y$ within $P$'s declared tolerances, for doses in $P$'s active region at
$t$.
\end{definition}

Three features of the definition do quiet work. First, the \emph{units
close}: alpha divided by alpha-per-activation yields activation, so the
prediction is dimensionally coherent by construction
\cite{buckingham}---which, after Conflict~\ref{ch04:conf:bridge}, is not a
triviality but a repair. The sources' inherited
$I_{\mathrm{pred}} = V/R_{\mathrm{total}}$ divided an intervention by a
score and produced dimensional nonsense; the candidate's version is
well-formed precisely because $r_{\mathrm{eff}}$ is required to be the
calibrated bridge quantity and nothing else. A constraint score in the
denominator is not a smaller error; it is a type error, rejected before
arithmetic.

Second, the \emph{strict positivity} of $r_{\mathrm{eff}}$ is a validation
rule, not a hope. A zero resistance would make the prediction undefined; a
negative one would smuggle a sign inversion into a quantity whose sign
conventions live elsewhere (in the direction's recorded sign). The
candidate's quantity validation rejects non-positive effective resistance
outright, so the prediction's denominator is well-defined whenever the
prediction is constructible at all.

Third, the prediction is \emph{local three times over}: to the profile (no
transfer across digests without a comparison rule), to the test point (no
transfer across layers, positions, directions, or hooks), and to the active
region (no claim at doses classified cutoff, saturated, breakdown,
non-ohmic, or unstable). The equation's resemblance to a global physical
law is the analogy the framework borrows; the typed reality is a patch of
claimed linearity, fenced on all sides.

\section{Admissibility Conditions}\label{ch06:admissibility}

The prediction may be \emph{evaluated} whenever its operands exist; it may
be \emph{admitted}---used as the reference against which residuals are
computed, regions confirmed, or hypotheses scored---only under conditions
that the implementation enforces as guards. Stated as a rule:

\begin{ruleenv}[Admissibility of the response prediction]\label{ch06:rule:admissibility}
A prediction $\hat y = a/r_{\mathrm{eff}}$ is admissible only if:
\begin{enumerate}
\item $r_{\mathrm{eff}}$ is strictly positive and carries an uncertainty
record produced by the profile's declared uncertainty method;
\item $r_{\mathrm{eff}}$ binds the profile's calibration and validation
split references---the data it was fit on and the data it was scored
on---as resolvable artifacts;
\item $r_{\mathrm{eff}}$ was produced by the profile's declared fit method,
with the fit receipt resolvable;
\item $r_{\mathrm{eff}}$ holds the role \texttt{CALIBRATED\_HELD\_OUT}, or
the role \texttt{NETWORK\_EQUIVALENT} with the separate validation that
role requires;
\item the dose $a$ lies within the profile's declared alpha range, and the
operating classification at the test point is \texttt{ACTIVE} on the
region axis;
\item the operation context is admitted and current: profile
\texttt{REVIEWED}, drift \texttt{CURRENT}, replay time within validity.
\end{enumerate}
\end{ruleenv}

The role vocabulary in condition (4) is the load-bearing novelty. The
candidate types every effective resistance with a \emph{resistance role},
an enumeration of four values, and the roles are not interchangeable:

\begin{description}
\item[\texttt{CALIBRATED\_HELD\_OUT}.] Fit on the calibration split,
scored on the disjoint validation split against the profile's tolerances.
The only role, besides a validated network equivalent, that admits
prediction. The role's constructor requirements enforce a held-out split
reference: a resistance cannot even claim the role without binding the
data that scored it.
\item[\texttt{APPARENT\_SAME\_SAMPLE}.] Computed as $a/y$ from the very
observations it would ``predict.'' The candidate allows this object to
exist---it is a legitimate descriptive statistic, the observed dose-response
ratio---but the role bars it from admission as a predictor. Its name is its
warning label.
\item[\texttt{NETWORK\_EQUIVALENT}.] Produced by a capability-gated
structural composition (series or parallel aggregation,
Chapter~\ref{ch09:networks} where present). Factory-constructed only: the
implementation refuses hand-built instances, so a network equivalent
carries its full derivation or does not exist. Admissible for prediction
only under the separate validation its gating demands.
\item[\texttt{DESIGN\_PROPOSAL}.] A declared value---someone's proposal for
what the resistance ought to be, typed as design information. Never
admissible for prediction; useful as a registered target that later
calibration may confirm or refute.
\end{description}

The enumeration is the anti-circularity rule made structural. In an untyped
system, the same floating-point number could serve as descriptive ratio,
fitted parameter, composed aggregate, and design guess, with its history
carried in comments if at all. Here the four histories are four types, the
guards read the type, and the tautological one is stopped at the door by
its name.

\section{The Anti-Circularity Rule}\label{ch06:anticirc}

\begin{ruleenv}[Anti-circularity]\label{ch06:rule:anticirc}
The following procedure is rejected regardless of its numerical outcome:
compute $r_{\mathrm{same}} = a/y$ on a sample of observations; compute
$\hat y = a / r_{\mathrm{same}}$ on that same sample; report the agreement
of $\hat y$ with $y$ as validation of the response model. A defensible test
fits $r_{\mathrm{eff}}$ on a registered calibration split and evaluates its
predictions on a disjoint held-out split, with uncertainty and discrepancy
represented \cite{kennedyohagan}.
\end{ruleenv}

The rejection is not a stylistic preference; the rejected procedure is
mathematically incapable of failing, and a test that cannot fail confirms
nothing.

\begin{proposition}[Same-sample identity is tautological]\label{ch06:prop:tautology}
Let $(a, y)$ be any observation with $a \neq 0$ and $y > 0$. Define
$r_{\mathrm{same}} = a / y$ and $\hat y = a / r_{\mathrm{same}}$. Then
$\hat y = y$ identically, for every such observation, regardless of whether
any linear (or any other) relation holds between dose and response.
\end{proposition}

\begin{proof}
By substitution:
$\hat y = a / r_{\mathrm{same}} = a / (a/y) = a \cdot (y/a) = y$. The
computation uses only the field axioms; no property of the data
enters. Hence the ``agreement'' $\hat y = y$ is an algebraic identity, true
for observations drawn from a linear system, a saturating system, a random
number generator, or a constant. Observing it carries zero bits of
information about the dose-response relation.
\end{proof}

\begin{remark}\label{ch06:rem:pointwise}
The multi-point version is only superficially better. Given $n$
observations $(a_i, y_i)$, computing per-point ratios
$r_i = a_i / y_i$ and ``predicting'' each $y_i$ from its own $r_i$ commits
the identity $n$ times. Fitting a single $r$ to all $n$ points and scoring
it on the same $n$ points breaks the exact identity but not the epistemic
problem: the score now measures in-sample goodness of fit, which improves
mechanically with every degree of freedom granted to the fit and says
nothing about behavior on data the fit has not seen. The rejected procedure
and its multi-point relatives sit on a continuum of self-scoring; the rule
draws the admissibility line where the scoring data and the fitting data
are disjoint by construction.
\end{remark}

\subsection{The statistical framing}\label{ch06:statistics}

Why is disjointness the line? The statistical answer has three layers.

\emph{Layer one: selection optimism.} A parameter chosen to minimize error
on a sample is, by that very selection, tuned to the sample's noise as well
as its signal. Its error on the selection sample is therefore biased
low---optimistically---relative to its error on fresh data, and the bias
grows with the flexibility of the fit relative to the data. This is the
elementary overfitting phenomenon, and the one-parameter case at issue here
is its minimal instance: even a single fitted slope captures some noise,
and a same-sample score inherits that optimism. Held-out evaluation removes
the bias by breaking the dependence between the noise the fit absorbed and
the noise the score encounters.

\emph{Layer two: the hypothesis is the shape, not the number.} What the
response model actually asserts is not ``there exists a number $r$ such
that $y \approx a/r$ on these points''---for any single point that is
trivially true, and for $n$ points it is true whenever the data are roughly
proportional---but ``the dose-response relation at this test point, in this
region, is stably linear, such that a slope estimated yesterday predicts
tomorrow.'' Stability under transfer from fitting data to fresh data is the
entire content of the claim. A same-sample test does not weakly test this
claim; it does not touch it. The train/test separation is thus not a
methodological garnish but the operational definition of the thing being
asserted.

\emph{Layer three: discrepancy is not noise.} Even a held-out test can
mislead if the model's error structure is misrepresented. The
Kennedy--O'Hagan calibration framework makes a distinction the candidate
adopts wholesale: parameter uncertainty (our estimate of
$r_{\mathrm{eff}}$ is imprecise) is not model discrepancy (the true
dose-response curve is not exactly linear, so even the best possible
$r_{\mathrm{eff}}$ leaves a structured, dose-dependent residual)
\cite{kennedyohagan}. A calibration that forces all misfit into parameter
uncertainty will report a confidently wrong slope with honest-looking error
bars; one that models discrepancy separately learns where the linear
hypothesis degrades---which, for a framework whose region classifier must
detect the edges of linearity, is not a nuisance term but the signal that
draws the fence. The admissibility conditions therefore require the fit
receipt to represent both, and the tolerance profile's linearity minimum is
the declared threshold at which discrepancy evidence disqualifies the
active-region claim itself.

\begin{remark}[Relation to the audit]\label{ch06:rem:audit}
Anti-circularity is the response model's local instance of the audit-wide
pattern of Chapter~\ref{ch04:audit}: a numerically computable operation
(divide, refit, compare) whose availability tempts users into an epistemic
claim its computation cannot support. The residual conflict
(Conflict~\ref{ch04:conf:residual}) documented the same pattern one step
downstream---residuals promoted to causes---and the network conflict
(Conflict~\ref{ch04:conf:networks}) one step up---compositions promoted to
laws. In all three cases the candidate's answer is identical: keep the
computation, type its output honestly, and gate the claim on evidence the
computation itself cannot manufacture.
\end{remark}

\section{Worked Hypothetical Illustrations}\label{ch06:illustrations}

The following examples use fabricated numbers chosen for arithmetic
clarity. \emph{They are illustrations of logic, not measurements: no model
was probed, no activation observed, and no value below describes any real
system.} Units are the candidate's model units (alpha, activation,
alpha-per-activation) and are as hypothetical as the numbers.

\begin{example}[The circular procedure cannot fail]\label{ch06:ex:circular}
\emph{(Hypothetical.)} Suppose a single sweep observation at dose
$a = 2.0$ alpha yields response $y = 0.8$ activation. The circular analyst
computes $r_{\mathrm{same}} = 2.0 / 0.8 = 2.5$ alpha-per-activation, then
``predicts'' $\hat y = 2.0 / 2.5 = 0.8$ activation, matching observation
exactly, and declares the relation validated.

Now corrupt the data deliberately: replace the response with a value drawn
from a table of random digits, say $y = 0.137$. The same procedure gives
$r_{\mathrm{same}} = 2.0 / 0.137 \approx 14.599$ and
$\hat y = 2.0 / 14.599 \approx 0.137$: perfect agreement again. The
procedure certifies random digits as obeying the response law with the
same confidence as any real signal---which is to say, its certification is
worthless (Proposition~\ref{ch06:prop:tautology}). In the candidate's
types, both computed ratios are constructible, but each is an
\texttt{EFFECTIVE\_RESISTANCE} with role
\texttt{APPARENT\_SAME\_SAMPLE}, and Rule~\ref{ch06:rule:admissibility}
refuses them admission as predictors. The refusal costs nothing real: the
descriptive ratio remains available as description.
\end{example}

\begin{example}[A held-out test that fails informatively]\label{ch06:ex:fail}
\emph{(Hypothetical.)} Suppose a calibration split contains sweep
observations at doses $(0.5, 1.0, 1.5)$ alpha with responses
$(0.25, 0.50, 0.75)$ activation---exactly proportional, slope
$0.5$ activation per alpha, so the declared fit method estimates
$r_{\mathrm{eff}} = 2.0$ alpha-per-activation with (say) a bootstrap
interval $[1.9, 2.1]$. In-sample agreement is perfect. The validation split
probes the withheld doses $(2.0, 2.5)$ and observes responses
$(0.82, 0.84)$ activation, whereas the model predicts
$\hat y = 2.0/2.0 = 1.00$ and $2.5/2.0 = 1.25$. Against a hypothetical
absolute tolerance of $0.05$ activation, both predictions miss badly
(errors $0.18$ and $0.41$).

The failure is informative in exactly the way the circular test could
never be: the response flattens near $0.8$--$0.85$ activation while
predictions keep climbing, the signature of entry into a saturation
region. The correct disposition under the candidate's machinery is not
``the model is refuted everywhere'' but ``the active region does not
extend to these doses'': the region classifier should bound the
active region below $2.0$ alpha, the resistance's claim contracts to the
bounded region, and the saturated observations are censored from residual
diagnostics per Conflict~\ref{ch04:conf:residual}'s rule. The held-out
design converted disagreement into a map of where the linear hypothesis
lives---precisely the discrepancy structure that Kennedy--O'Hagan-style
calibration is built to expose \cite{kennedyohagan}.
\end{example}

\begin{example}[A held-out test that passes, and what it would
license]\label{ch06:ex:pass}
\emph{(Hypothetical.)} Suppose instead the validation observations at
doses $(0.75, 1.25)$ alpha---interior to the calibrated range---are
$(0.36, 0.64)$ activation against predictions $(0.375, 0.625)$: errors of
$0.015$ within the hypothetical absolute tolerance $0.05$, with the
linearity minimum met and off-target contamination below its ceiling.
Then, and only then, may the fitted quantity hold role
\texttt{CALIBRATED\_HELD\_OUT}, and predictions become admissible---under
this profile, at this test point, for doses inside the calibrated range,
in the active region, while the profile remains current.

Note how little even the pass licenses. It does not extend one dose beyond
the grid's endpoints (extrapolation was never scored). It does not
transfer to the neighboring layer, a re-derived direction, or next month's
checkpoint (different test point or profile). It does not assert that the
intervention \emph{causes} the response in the interventionist sense
without the causal-model apparatus that claim requires \cite{pearl}. And
it is two hypothetical numbers agreeing with two other hypothetical
numbers: in the real program, the pass threshold is a registered
hypothesis with controls (zero-dose, random-direction, wrong-layer,
label-shuffle, norm-matched) whose defeat is also required---agreement
alone, without superiority over matched controls, is not evidence of a
direction-specific response at all.
\end{example}

\section{Verification Versus Validation}\label{ch06:vandv}

Sargent's classical decomposition distinguishes the questions a simulation
or model effort must keep separate: conceptual model validation (is the
model's theory a defensible account of the problem entity?), computerized
model verification (does the implementation correctly implement the
specification?), and operational validation (do the implemented model's
outputs agree with the real system's behavior in the domain of intended
use?) \cite{sargent}. The candidate maps onto the decomposition cleanly,
and the mapping is the cleanest way to state what has and has not been
done.

\emph{Conceptual model validation}: partially addressed, negatively. The
source audit (Chapter~\ref{ch04:audit}) establishes that the inherited
conceptual model was not even internally univocal---eight conflicts---and
the candidate's typed reconstruction restores internal coherence. Internal
coherence is the precondition of conceptual validity, not its
establishment; whether a local linear dose-response theory is a
\emph{defensible account} of steered activations is exactly what the
registered hypotheses will assess.

\emph{Computerized model verification}: this is what the bounded software
evidence actually is. The test suite at the Gate~4 receipt (65 passing test
methods, 58 parameterized subtests) exercises the guards, the types, the
canonical encoding, the split-distinctness rules, the exact quantization
boundaries, the fail-closed behaviors---the implementation against its
specification. Synthetic known-fault boards (linear, cutoff, saturating,
nonlinear, parallel, confounded, bypass) verify that the classifier and
guards respond to designed inputs as specified. All of this is
verification in Sargent's sense: it establishes that the software does what
the design says, on inputs the designers constructed. It establishes
nothing about any real model's behavior.

\emph{Operational validation}: not performed. No live system has been
calibrated; no held-out prediction against a real activation has been
scored; the registered hypotheses are open. The candidate's entire status
vocabulary---\textsc{ring 4 / not licensed / held}, constant-false
authenticated eligibility, roles that no real quantity yet holds---exists
to keep this cell of the matrix visibly empty until the preregistered
closure protocol fills it.

\begin{remark}\label{ch06:rem:vandv}
The distinction dissolves a misreading to which test counts invite
themselves. ``65 tests pass'' is evidence about the correspondence between
code and specification; it is categorically incapable of being evidence
about the correspondence between specification and world. A perfectly
verified implementation of a false theory passes all its tests. Conversely,
operational validation without verification is uninterpretable---an
agreement between a buggy implementation and reality confirms nothing about
the theory. The program needs both, in order, and it currently has the
first only.
\end{remark}

\section{What a Validated Response Model Would and Would Not
License}\label{ch06:license}

Suppose---counterfactually as of this writing---that the closure protocol
ran to completion for some profile: preregistered claim, frozen
environment, randomized held-out sweeps, all controls defeated, independent
recomputation, replication on a held-out checkpoint under a prespecified
comparison rule, independent review. What then?

\emph{It would license}: interpolative prediction of directional activation
response at the calibrated test points, within the calibrated dose range,
in the active region, under the current profile, with stated uncertainty;
the use of such predictions as reference curves for residual
\emph{ranking} (never identification) per
Conflict~\ref{ch04:conf:residual}; region classification with the
confidence the validated fit supports; and---if the corresponding
composition hypotheses also closed---the network equivalences within their
own validated scope. It would, in short, license the framework's
\emph{instrumental} uses: prediction, screening, prioritization, all
profile-local.

\emph{It would not license}: causal claims (``the dose causes the response
via mechanism $M$'')---interventionist causal language requires an explicit
causal model and identifiable interventions, which a fitted slope is not
\cite{pearl}; physical claims of any kind---the units remain model units,
and the forbidden-alias rule keeps volt/ampere/ohm out of reach
permanently, validation or no; transfer---to other profiles, models,
layers, directions, corpora, or dtypes---without separately admitted
comparison rules; aggregate or league-table claims across systems;
``measures attention'' or any construct-identification claim, which would
require its own validity argument connecting the projection to the
construct; safety verdicts, product capabilities, or marketing predicates,
which are licensing acts reserved to authorities the candidate does not
contain; and permanence---a validated profile drifts, expires, and returns
to \texttt{UNASSESSED} obligations like any other.

The asymmetry of the two lists is the chapter's summary. Validation, even
at its best, buys a bounded instrument: a locally trustworthy slope with
error bars, inside a fence. Everything outside the fence---cause, physics,
transfer, safety, authority---requires evidence of a different kind, which
no amount of dose-response agreement accumulates. The candidate's types
are built so that the fence travels with every number; the purpose of this
chapter has been to show that the fence is not caution appended to the
science but the exact shape of what the science, done honestly, would have
shown.

\begin{boundary}\label{ch06:bound}
This chapter defines a hypothesis and its admission conditions. The
response model is unvalidated: no admissible \texttt{CALIBRATED\_HELD\_OUT}
resistance exists for any live system, all worked numbers above are
fabricated illustrations, and nothing here establishes empirical accuracy,
causal interpretation, physical identification, or licensed use. The
anti-circularity rule and admissibility guards are implemented and
verification-tested; their scientific worth is exactly the worth of the
experiments they are waiting for.
\end{boundary}
